\documentclass[conference]{IEEEtran}
\IEEEoverridecommandlockouts
\usepackage{cite}
\usepackage{amsmath}
\usepackage{graphicx}
\usepackage{booktabs}

% Tighten spacing globally
\setlength{\abovecaptionskip}{3pt}
\setlength{\belowcaptionskip}{0pt}
\setlength{\floatsep}{4pt}
\setlength{\textfloatsep}{6pt}
\setlength{\intextsep}{4pt}

\begin{document}

\title{QLoRA Fine-Tuning of PaliGemma for Remote Sensing Image Captioning}

\author{
\IEEEauthorblockN{\.{I}layda Erg\"{u}r}
\IEEEauthorblockA{\textit{DI725 --- Transformers and Attention-Based Deep Networks}\\
Middle East Technical University, Ankara, Turkey}
}

\maketitle

\begin{abstract}
We apply QLoRA to fine-tune PaliGemma-3B on the RISC remote sensing image captioning dataset,
training only 0.39\% of model parameters.
Fine-tuning yields substantial gains over the pre-trained baseline:
BLEU-4 improves from 0.0062 to 0.0553, METEOR from 0.0897 to 0.2285,
and ROUGE-L from 0.1476 to 0.2675.
\end{abstract}

\section{Introduction}

PaliGemma~\cite{beyer2024paligemma} is a 3B-parameter VLM combining the SigLIP vision encoder
with the Gemma language decoder.
While capable in general domains, it underperforms on specialised tasks without adaptation.
Full fine-tuning is memory-prohibitive, so we apply QLoRA~\cite{dettmers2023qlora}
--- 4-bit quantisation combined with low-rank adapter weights ---
to efficiently adapt PaliGemma to remote sensing image captioning.

\section{Method}

\subsection{Dataset}
We use the RISC dataset~\cite{riscdata} (\texttt{caglarmert/full\_riscm}):
44,521 satellite images (224$\times$224\,px) with five captions each,
split 95/5 into 42,294 training pairs and 2,226 test images.

\subsection{Model and QLoRA Configuration}
PaliGemma-3B is loaded in 4-bit NF4 quantisation via BitsAndBytes.
LoRA adapters ($r\!=\!8$) are attached to all seven linear projection layers of the Gemma decoder
(\texttt{q, k, v, o, gate, up, down\_proj}), yielding \textbf{11.3M / 2.93B trainable parameters (0.39\%)}.
The SigLIP encoder is frozen throughout.
Training uses AdamW-8bit, learning rate $2\!\times\!10^{-5}$,
effective batch size 16, and 1 epoch over the full training set.

\section{Results}

\begin{table}[h]
\centering
\caption{Captioning metrics on 200 test images (single GT reference)}
\label{tab:results}
\setlength{\tabcolsep}{5pt}
\begin{tabular}{lccc}
\toprule
\textbf{Model} & \textbf{BLEU-4} & \textbf{METEOR} & \textbf{ROUGE-L} \\
\midrule
Baseline PaliGemma & 0.0062 & 0.0897 & 0.1476 \\
QLoRA Fine-tuned   & \textbf{0.0553} & \textbf{0.2285} & \textbf{0.2675} \\
\bottomrule
\end{tabular}
\end{table}

Table~\ref{tab:results} and Fig.~\ref{fig:metrics} show consistent improvement across all metrics.
BLEU-4 increases nearly 9$\times$; METEOR and ROUGE-L roughly double.
Low absolute BLEU-4 values are expected for single-reference evaluation.
Fig.~\ref{fig:examples} shows qualitative comparisons: the baseline generates generic descriptions,
while the fine-tuned model produces domain-specific captions
(e.g., runways, vegetation patterns, road networks).

\begin{figure}[h]
\centering
\includegraphics[width=\columnwidth, height=3.2cm, keepaspectratio]%
  {../figures/metrics_comparison.png}
\caption{BLEU-4, METEOR, ROUGE-L before and after QLoRA fine-tuning.}
\label{fig:metrics}
\end{figure}

\begin{figure}[h]
\centering
\includegraphics[width=\columnwidth, height=3.8cm, keepaspectratio]%
  {../figures/results_visualization.png}
\caption{Ground truth, baseline, and fine-tuned captions for five test images.}
\label{fig:examples}
\end{figure}

\section{Conclusion}
QLoRA enables effective domain adaptation of PaliGemma-3B with only 0.39\% trainable parameters,
yielding significant quantitative and qualitative improvements on remote sensing image captioning.

\begin{thebibliography}{1}

\bibitem{beyer2024paligemma}
L.~Beyer et al., ``PaliGemma: A versatile 3B VLM for transfer,''
\textit{arXiv:2407.07726}, 2024.

\bibitem{dettmers2023qlora}
T.~Dettmers et al., ``QLoRA: Efficient finetuning of quantized LLMs,''
\textit{NeurIPS}, 2023.

\bibitem{riscdata}
Y.~Qu et al., ``RISC: Remote sensing image captioning dataset,''
\textit{HuggingFace}, \texttt{caglarmert/full\_riscm}, 2024.

\end{thebibliography}

\end{document}
